\chapter{Simetri Sumbu}\label{ch:g6:symmetry}

Lipatlah selembar kertas sepanjang sebuah garis: setiap titik mendarat
di titik yang lain, yaitu bayangan cerminnya. Lipatan ini adalah
\emph{\omterm{def:g6:symmetry:def}{simetri sumbu}} (atau pencerminan), transformasi geometri yang
pertama dalam buku ini --- yang kedua, simetri pusat, menyusul pada
\cref{ch:g7:central}.

\section{Pencerminan}

\begin{definition}[Titik bayangan]\label{def:g6:symmetry:def}
Misalkan $d$ sebuah garis. \emph{Bayangan}\index{simetri sumbu} titik
$M$ terhadap $d$ adalah titik $M'$ sedemikian sehingga $d$ menjadi
\omterm{def:g6:symmetry:bisector}{garis sumbu} $[MM']$: \omterm{def:g6:lines:objects}{ruas garis} $[MM']$ \omterm{def:g6:lines:perp}{tegak lurus} terhadap $d$, dan
$d$ memotongnya di titik tengahnya. Titik yang terletak \emph{pada}
$d$ adalah bayangannya sendiri.
\end{definition}

\begin{omfigure}
\begin{tikzpicture}[scale=1.05]
  \draw[thick, omExo] (0,-1.7) -- (0,1.9) node[above, font=\small] {$d$};
  \fill (-2.1,1.0) circle (1.6pt) node[left, font=\small] {$M$};
  \fill (2.1,1.0) circle (1.6pt) node[right, font=\small] {$M'$};
  \draw[dashed] (-2.1,1.0) -- (2.1,1.0);
  \draw (-0.10,0.90) -- (0.10,1.10);
  \draw (-1.15,0.90) -- (-0.95,1.10);
  \draw (0.95,0.90) -- (1.15,1.10);
  \draw[thin] (0,1.0) ++(0.16,0) -- ++(0,0.16) -- ++(-0.16,0);
  \fill (-1.3,-1.1) circle (1.6pt) node[left, font=\small] {$N$};
  \fill (1.3,-1.1) circle (1.6pt) node[right, font=\small] {$N'$};
  \draw[dashed] (-1.3,-1.1) -- (1.3,-1.1);
  \fill (0,0.2) circle (1.6pt) node[above right, font=\small] {$P = P'$};
\end{tikzpicture}

{\small Pencerminan terhadap garis $d$: sumbunya \omterm{def:g6:lines:perp}{tegak lurus} terhadap
$[MM']$ dan melalui titik tengahnya (tanda garis yang sama). Titik $P$
pada sumbunya tidak berpindah.}
\end{omfigure}

\begin{method}[Melukis titik bayangan]\label{met:g6:symmetry:construct}
Untuk melukis bayangan $M'$ dari titik $M$ terhadap garis $d$:
\begin{enumerate}
  \item gambarlah \omterm{def:g6:lines:perp}{garis tegak lurus} terhadap $d$ melalui $M$ (dengan
        penggaris siku); sebutlah $H$ titik potongnya dengan $d$;
  \item ukurlah $MH$ dengan jangka atau penggaris;
  \item letakkan $M'$ pada \omterm{def:g6:lines:perp}{garis tegak lurus} yang sama, di sisi lain
        $d$, dengan $HM' = MH$.
\end{enumerate}
Di kertas berpetak caranya lebih cepat: bilanglah petak dari $M$ ke
sumbunya, lalu bilanglah sebanyak itu juga di sisi yang lain.
\end{method}

\begin{example}[Di kertas berpetak]\label{ex:g6:symmetry:grid}
Kalau sumbunya garis petak yang tegak dan $M$ berada $3$ petak di
sebelah kirinya, maka $M'$ berada $3$ petak di sebelah kanannya, pada
baris yang sama. Untuk sumbu yang miring, pembilangannya berjalan
sepanjang diagonal --- selalu \omterm{def:g6:lines:perp}{tegak lurus} terhadap sumbunya.
\end{example}

\begin{omfigure}
\begin{tikzpicture}[scale=0.55]
  \draw[gray!40, thin] (0,0) grid (12,6);
  \draw[thick, omExo] (6,-0.4) -- (6,6.4);
  \draw[thick, omDef, fill=omDef!15]
    (1,1) -- (4,1) -- (4,3) -- (3,3) -- (3,5) -- (1,5) -- cycle;
  \draw[thick, omThm, fill=omThm!15]
    (11,1) -- (8,1) -- (8,3) -- (9,3) -- (9,5) -- (11,5) -- cycle;
  \node[omDef, font=\small] at (2.4,0.45) {gambarnya};
  \node[omThm, font=\small] at (9.6,0.45) {bayangan cerminnya};
\end{tikzpicture}

{\small Sebuah gambar dan pencerminannya terhadap sumbu tegak: setiap
\omterm{def:g5:solids:def}{titik sudut} dicerminkan petak demi petak, dan kiri serta kanannya bertukar.}
\end{omfigure}

\begin{proposition}[Apa yang dipertahankan pencerminan]\label{prop:g6:symmetry:preserved}
\omterm{def:g6:symmetry:def}{Bayangan} sebuah gambar terhadap sebuah garis adalah gambar dengan
\emph{bentuk dan ukuran yang sama}: panjang, sudut, \omterm{def:g4:measure:perimeter}{keliling} dan luas
semuanya dipertahankan. \omterm{def:g6:symmetry:def}{Bayangan} \omterm{def:g6:lines:objects}{ruas garis} adalah \omterm{def:g6:lines:objects}{ruas garis},
bayangan lingkaran adalah lingkaran dengan \omterm{def:g3:shapes:shapes}{jari-jari} yang sama, dan
bayangan garis adalah garis.
\end{proposition}
\admitted

\begin{remark}
Satu hal \emph{tidak} dipertahankan: arah hadapnya. Huruf ``b'' yang
dicerminkan menjadi ``d'' --- tulisan cermin. Kalau menjiplak gambarmu
lalu membalik kertas kalkirnya menghasilkan bayangan itu,
pencerminanmu benar.
\end{remark}

\section{Sumbu simetri sebuah gambar}

\begin{definition}[Sumbu simetri]\label{def:g6:symmetry:axis}
Sebuah garis $d$ disebut \emph{sumbu simetri}\index{sumbu simetri}
sebuah gambar jika pencerminan terhadap $d$ mengirim gambar itu tepat
ke dirinya sendiri --- melipat sepanjang $d$ membuat kedua paruhnya berimpit.
\end{definition}

\begin{omfigure}
\begin{tikzpicture}[scale=0.62]
  % segitiga sama kaki
  \draw[thick] (0,0) -- (2.4,0) -- (1.2,2.6) -- cycle;
  \draw[omThm, dashed] (1.2,-0.5) -- (1.2,3.0);
  \node[below, font=\small] at (1.2,-0.7) {sama kaki: $1$ sumbu};
  % persegi panjang
  \begin{scope}[xshift=5.6cm]
  \draw[thick] (0,0.3) rectangle (3.2,2.3);
  \draw[omThm, dashed] (1.6,-0.4) -- (1.6,3.0);
  \draw[omThm, dashed] (-0.5,1.3) -- (3.7,1.3);
  \node[below, font=\small] at (1.6,-0.7) {persegi panjang: $2$ sumbu};
  \end{scope}
  % persegi
  \begin{scope}[xshift=12.0cm]
  \draw[thick] (0.3,0.1) rectangle (2.7,2.5);
  \draw[omThm, dashed] (1.5,-0.5) -- (1.5,3.1);
  \draw[omThm, dashed] (-0.3,1.3) -- (3.3,1.3);
  \draw[omThm, dashed] (0.0,-0.2) -- (3.0,2.8);
  \draw[omThm, dashed] (0.0,2.8) -- (3.0,-0.2);
  \node[below, font=\small] at (1.5,-0.7) {persegi: $4$ sumbu};
  \end{scope}
\end{tikzpicture}

{\small \omterm{def:g6:symmetry:axis}{Sumbu simetri} bangun yang biasa dikenal. Hati-hati dengan
persegi panjang: diagonalnya \emph{bukan} \omterm{def:g6:symmetry:axis}{sumbu simetri} (lipatlah
sepanjang diagonalnya --- kedua paruhnya tidak berimpit kecuali kalau ia persegi).}
\end{omfigure}

\begin{example}\label{ex:g6:symmetry:shapes}
Membilang \omterm{def:g6:symmetry:axis}{sumbu simetri}:
\begin{itemize}
  \item \omterm{def:g6:shapes:triangles}{segitiga sama kaki}: $1$ (melalui \omterm{def:g5:solids:def}{titik sudut} puncak dan
        tengah alasnya); \omterm{def:g6:shapes:triangles}{segitiga sama sisi}: $3$;
  \item persegi panjang: $2$ (kedua ``garis tengahnya'', \emph{bukan}
        diagonalnya); \omterm{def:g6:shapes:quadrilaterals}{belah ketupat}: $2$ (yaitu diagonalnya!); persegi: $4$;
  \item lingkaran: setiap garis yang melalui pusatnya --- tak hingga banyaknya.
\end{itemize}
\end{example}

\begin{example}[Simetri membuktikan kesamaan]\label{ex:g6:symmetry:proves}
Mengapa kedua sudut alas \omterm{def:g6:shapes:triangles}{segitiga sama kaki} sama besar? Lipatlah
segitiga itu sepanjang \omterm{def:g6:symmetry:axis}{sumbu simetrinya}: paruh kirinya mendarat tepat
pada paruh kanannya, jadi sudut alas kiri mendarat pada sudut alas
kanan --- keduanya pasti berukuran sama
(\cref{prop:g6:symmetry:preserved}: pencerminan mempertahankan sudut).
\end{example}

\section{Garis sumbu}

\begin{definition}[Garis sumbu]\label{def:g6:symmetry:bisector}
\emph{Garis sumbu}\index{garis sumbu} sebuah \omterm{def:g6:lines:objects}{ruas garis} $[AB]$ adalah
garis yang \omterm{def:g6:lines:perp}{tegak lurus} terhadap $[AB]$ dan melalui titik tengahnya ---
ia adalah \omterm{def:g6:symmetry:axis}{sumbu simetri} yang mempertukarkan $A$ dan $B$.
\end{definition}

\begin{proposition}[Berjarak sama]\label{prop:g6:symmetry:bisector}
Sebuah titik $M$ terletak pada \omterm{def:g6:symmetry:bisector}{garis sumbu} $[AB]$ tepat ketika ia
berjarak sama dari $A$ dan dari $B$: $MA = MB$.
\end{proposition}
\admitted

\begin{method}[Lukisan dengan jangka]\label{met:g6:symmetry:bisector}
Untuk menggambar \omterm{def:g6:symmetry:bisector}{garis sumbu} $[AB]$ tanpa mengukur:
\begin{enumerate}
  \item bukalah jangka lebih dari setengah $AB$;
  \item gambarlah dua busur berpusat di $A$, di atas dan di bawah ruas
        garisnya;
  \item dengan bukaan yang \emph{sama}, gambarlah dua busur berpusat
        di $B$;
  \item hubungkan kedua titik potongnya: garis inilah \omterm{def:g6:symmetry:bisector}{garis sumbunya}
        (setiap titik potong berjarak sama dari $A$ dan $B$).
\end{enumerate}
\end{method}

\section{Latihan}

\begin{exercise}[$\star$]\label{exo:g6:symmetry:1}
Salinlah di kertas berpetak: sumbu tegak $d$, titik $A$ dua petak di
sebelah kirinya, titik $B$ lima petak di sebelah kirinya dan satu
petak lebih tinggi. Lukislah titik bayangan $A'$ dan $B'$ terhadap $d$.
\end{exercise}

\begin{exercise}[$\star$]\label{exo:g6:symmetry:2}
Gambarlah sebuah garis $d$ dan titik $M$ berjarak $3$ cm dari $d$
(bukan di kertas berpetak). Lukislah bayangan $M'$ dari $M$ terhadap
$d$ dengan penggaris siku dan penggaris, mengikuti
\cref{met:g6:symmetry:construct}. Berapa jarak antara $M$ dan $M'$?
\end{exercise}

\begin{exercise}[$\star$]\label{exo:g6:symmetry:3}
Huruf kapital mana dalam abjad, yang ditulis dalam bentuknya yang
paling sederhana, yang punya \omterm{def:g6:symmetry:axis}{sumbu simetri} tegak (seperti A)? Yang
mendatar (seperti B)? Keduanya?
\end{exercise}

\begin{exercise}[$\star$]\label{exo:g6:symmetry:4}
Berapa \omterm{def:g6:symmetry:axis}{sumbu simetri} yang dipunyai: \omterm{def:g6:shapes:triangles}{segitiga sama sisi}? persegi
panjang (yang bukan persegi)? \omterm{def:g6:shapes:quadrilaterals}{belah ketupat} (yang bukan persegi)?
persegi? lingkaran? Gambarlah setiap bangun beserta sumbunya.
\end{exercise}

\begin{exercise}[$\star$]\label{exo:g6:symmetry:5}
Benar atau salah? ``Diagonal sebuah persegi panjang adalah \omterm{def:g6:symmetry:axis}{sumbu
simetri} persegi panjang itu.'' Jelaskan dengan alasan lipatan atau dengan gambar.
\end{exercise}

\begin{exercise}[$\star$]\label{exo:g6:symmetry:6}
Segitiga $ABC$ dicerminkan terhadap garis $d$ menjadi $A'B'C'$.
Diketahui $AB = 4$ cm, $\widehat{BAC} = 70^\circ$ dan \omterm{def:g4:measure:perimeter}{keliling} $ABC$
adalah $12$ cm: sebutkan, tanpa melukis apa pun, $A'B'$,
$\widehat{B'A'C'}$ dan \omterm{def:g4:measure:perimeter}{keliling} $A'B'C'$.
\end{exercise}

\begin{exercise}[$\star$]\label{exo:g6:symmetry:7}
Gambarlah \omterm{def:g6:lines:objects}{ruas garis} $[AB]$ sepanjang $7$ cm lalu lukislah \omterm{def:g6:symmetry:bisector}{garis
sumbunya} dengan jangka (\cref{met:g6:symmetry:bisector}). Pilihlah
sembarang titik $M$ padanya lalu periksalah dengan penggaris bahwa $MA = MB$.
\end{exercise}

\begin{exercise}[$\star\star$]\label{exo:g6:symmetry:8}
Dua pohon berdiri di titik $A$ dan $B$ sebuah kebun yang datar. Di
mana air mancur dapat ditanam agar berjarak sama dari kedua pohon itu?
Jelaskan \emph{semua} tempat yang mungkin.
\end{exercise}

\begin{exercise}[$\star\star$]\label{exo:g6:symmetry:9}
Di kertas berpetak, gambarlah sumbu $d$ (diagonal petak pada
$45^\circ$) dan sebuah gambar kecil berbentuk L di salah satu sisinya.
Lukislah gambar bayangannya terhadap $d$. (Cerminkan setiap \omterm{def:g5:solids:def}{titik
sudut}: untuk sumbu $45^\circ$, titik yang berjarak $2$ petak di kanan
sumbu mendarat $2$ petak di atasnya.)
\end{exercise}

\begin{exercise}[$\star\star$]\label{exo:g6:symmetry:10}
Gambarlah lingkaran berpusat $O$, garis $d$ yang melalui $O$, dan
titik $M$ pada lingkaran. Di mana bayangan $M$ terhadap $d$? Jelaskan
mengapa pencerminan itu mengirim lingkarannya ke dirinya sendiri.
\end{exercise}

\begin{exercise}[$\star\star$]\label{exo:g6:symmetry:11}
Dengan memakai sifat berjarak sama
(\cref{prop:g6:symmetry:bisector}), jelaskan mengapa kedua titik
potong busur pada \cref{met:g6:symmetry:bisector} benar-benar terletak
pada \omterm{def:g6:symmetry:bisector}{garis sumbu} $[AB]$.
\end{exercise}

\begin{exercise}[$\star\star\star$]\label{exo:g6:symmetry:12}
Sebuah bola bilyar di titik $A$ harus memantul pada dinding lurus $d$
lalu mencapai titik $B$ (keduanya di sisi dinding yang sama).
Cerminkan $B$ terhadap $d$ menjadi $B'$, lalu gambarlah \omterm{def:g6:lines:objects}{ruas garis}
$[AB']$. Jelaskan mengapa titik pantul yang terbaik adalah tempat
$[AB']$ memotong dindingnya. (Gagasannya: jalur $A \to P \to B$ sama
panjang dengan $A \to P \to B'$.)
\end{exercise}

\section{Soal: Air mancur, piring yang pecah, dan cermin yang
berjalan}

\begin{problem}[{Soal akhir pekan --- ketiga garis sumbu sebuah
segitiga bertemu di satu titik: lingkaran melalui tiga titik, dan apa
yang dikerjakan dua cermin bersama-sama}]
\label{pb:g6:symmetry:1}
Pada \cref{exo:g6:symmetry:8} sebuah air mancur harus berdiri
berjarak sama dari \emph{dua} pohon, dan jawabannya adalah satu
garis penuh berisi tempat. Soal ini menanam pohon yang \emph{ketiga}
--- dan seluruh garis itu mengerut menjadi satu titik yang tertentu
sempurna. Titik itu menyembunyikan harta karun geometri: ia
memungkinkanmu menggambar satu-satunya lingkaran yang melalui tiga
titik yang diberikan, membangun kembali seluruh piring dari sekeping
\omterm{def:g6:fractions:def}{pecahan}, dan ia dimiliki setiap segitiga. Bagian terakhirnya
melakukan percobaan dengan dua cermin yang hasilnya mengumumkan
geometri tahun depan.

\medskip
\noindent\textbf{Bagian I --- Dua pohon, lalu peta wilayah.}
\begin{enumerate}
  \item Gambarlah dua titik $A$ dan $B$ dengan $AB = 6$ cm lalu
        lukislah \omterm{def:g6:symmetry:bisector}{garis sumbu} $[AB]$ dengan jangka
        (\cref{met:g6:symmetry:bisector}). Pilihlah sebuah titik $M$
        padanya lalu periksalah dengan penggaris bahwa $MA = MB$.
  \item \Cref{prop:g6:symmetry:bisector} mengatakan dua hal
        sekaligus: setiap titik \emph{pada} \omterm{def:g6:symmetry:bisector}{garis sumbunya} berjarak
        sama dari $A$ dan $B$, dan setiap titik yang berjarak sama
        \emph{terletak pada} \omterm{def:g6:symmetry:bisector}{garis sumbu} itu. Yang mana di antara
        keduanya yang menjawab pertanyaan air mancur pada
        \cref{exo:g6:symmetry:8}? Apa yang dikatakan proposisi itu
        tentang titik yang \emph{tidak} pada \omterm{def:g6:symmetry:bisector}{garis sumbunya}?
  \item Ambillah titik $M$ di sisi \omterm{def:g6:symmetry:bisector}{garis sumbu} yang sama dengan $A$.
        \omterm{def:g6:lines:objects}{Ruas garis} $[MB]$ memotong \omterm{def:g6:symmetry:bisector}{garis sumbu} itu di titik $P$.
        Bandingkan jalur $M \to P \to B$ dengan jalur
        $M \to P \to A$, lalu jelaskan mengapa $M$ lebih dekat ke
        $A$ daripada ke $B$: \omterm{def:g6:symmetry:bisector}{garis sumbunya} membelah kertas menjadi
        ``sisi $A$'' dan ``sisi $B$''.
  \item Dua sekolah $A$ dan $B$ melayani sebuah kota. Setiap anak
        berjalan ke sekolah yang lebih dekat. Gambarlah peta kecil
        dengan $A$ dan $B$, lalu warnailah daerah kota yang tepat
        dilayani sekolah $A$. Garis apa yang menjadi batasnya?
  \item Sekolah ketiga $C$ dibuka (letakkan agar ketiga sekolahnya
        membentuk segitiga). Lukislah \omterm{def:g6:symmetry:bisector}{garis sumbu} $[AB]$ dan \omterm{def:g6:symmetry:bisector}{garis
        sumbu} $[BC]$, lalu sebutlah $O$ titik potongnya. Dua kesamaan
        jarak apa yang dipenuhi $O$, dan mengapa?
\end{enumerate}

\medskip
\noindent\textbf{Bagian II --- Tiga pohon, satu titik, satu
lingkaran.}
\begin{enumerate}[resume]
  \item Simpulkan dari pertanyaan 5 bahwa $OA = OC$ --- dan karena
        itu $O$ juga terletak pada \omterm{def:g6:symmetry:bisector}{garis sumbu} $[AC]$, yang tidak
        pernah kamu gambar. Simpulkan: \emph{ketiga \omterm{def:g6:symmetry:bisector}{garis sumbu} sisi
        sebuah segitiga semuanya melalui satu titik.}
  \item Jelaskan mengapa lingkaran berpusat $O$ dengan \omterm{def:g3:shapes:shapes}{jari-jari} $OA$
        juga melalui $B$ dan melalui $C$: itulah \emph{satu-satunya}
        lingkaran yang melalui ketiga titik itu --- \emph{lingkaran
        luar}\index{lingkaran luar} segitiga $ABC$.
  \item Lukislah segitiga dengan sisi $6$ cm, $5$ cm dan $4$ cm
        (\cref{met:g6:shapes:construct}), lalu \omterm{pb:g6:symmetry:1}{lingkaran luarnya}.
        Mengapa cukup menggambar \emph{dua} \omterm{def:g6:symmetry:bisector}{garis sumbu} saja untuk
        menemukan pusatnya?
  \item Piring yang pecah: seorang arkeolog menggali sekeping \omterm{def:g6:fractions:def}{pecahan}
        yang bagian utuhnya hanya sepotong tepi lingkaran piring itu.
        Jelaskan resep yang memulihkan pusat dan \omterm{def:g3:shapes:shapes}{jari-jari} piring
        aslinya --- lalu ujilah: gambarlah busur lingkaran dengan
        jangkamu, ``lupakan'' pusatnya, tandailah tiga titik pada
        busurnya, lalu bangunlah kembali.
  \item Mungkinkah tiga titik sama sekali \emph{tidak} punya
        lingkaran yang melaluinya? Andaikan $A$, $B$, $C$ segaris.
        Apa yang dapat dikatakan tentang \omterm{def:g6:symmetry:bisector}{garis sumbu} $[AB]$ dan \omterm{def:g6:symmetry:bisector}{garis
        sumbu} $[BC]$ (\cref{prop:g6:lines:facts})? Simpulkan.
\end{enumerate}

\medskip
\noindent\textbf{Bagian III --- Sumbu di mana-mana, dan cermin yang
berjalan.}
\begin{enumerate}[resume]
  \item \Cref{ex:g6:symmetry:shapes} membilang $3$ sumbu untuk
        \omterm{def:g6:shapes:triangles}{segitiga sama sisi} dan $4$ untuk persegi. Terkalah
        hitungannya untuk \omterm{def:g4:geometry:polygon}{segi lima} beraturan dan \omterm{def:g4:geometry:polygon}{segi enam}
        beraturan, lalu jelaskan di mana sumbunya lewat pada tiap
        keadaan (banyak sisi yang \omterm{def:g3:numbers:evenodd}{ganjil} melawan yang \omterm{def:g3:numbers:evenodd}{genap}: kedua
        keadaannya berbeda).
  \item Lipatlah selembar kertas sekali lalu guntinglah sebuah bentuk
        menembus kedua lapisannya; bukalah lipatannya. Mengapa
        lubangnya \emph{selalu} punya \omterm{def:g6:symmetry:axis}{sumbu simetri}, apa pun yang
        kamu gunting, dan garis mana yang menjadi sumbunya?
  \item Lipatlah kertasnya dua kali --- lipatan kedua \omterm{def:g6:lines:perp}{tegak lurus}
        terhadap yang pertama --- guntinglah, lalu bukalah. Berapa
        \omterm{def:g6:symmetry:axis}{sumbu simetri} yang paling sedikit dipunyai lubangnya, dan
        berapa salinan guntinganmu yang kamu lihat? (Kepingan salju
        kertas adalah gagasan ini yang dilipat lebih jauh.)
  \item Cermin yang berjalan: di kertas berpetak, gambarlah sumbu
        tegak $d_1$, sebuah bendera kecil $F$ di sebelah kirinya, dan
        sumbu tegak kedua $d_2$, $4$ petak di sebelah kanan $d_1$.
        Cerminkan $F$ terhadap $d_1$ (bayangannya $F'$), lalu
        cerminkan $F'$ terhadap $d_2$ (bayangannya $F''$).
        Bandingkan $F$ dan $F''$: arah hadapnya sama atau tercermin?
        Berapa petak gambarnya berpindah? Bandingkan dengan jarak
        antara kedua sumbunya.
  \item Percobaan yang sama dengan $d_2$ \emph{\omterm{def:g6:lines:perp}{tegak lurus}} terhadap
        $d_1$: cerminkan benderanya terhadap sumbu tegak, lalu
        bayangannya terhadap sumbu mendatar. Jelaskan letak $F''$
        terhadap $F$ dan terhadap titik potong kedua sumbunya. (Kamu
        baru saja menemukan transformasi tahun depan:
        \cref{ch:g7:central}.)

\end{enumerate}
\end{problem}
